\begin{lemma}[The penultimate full-tail prefix is elementary]
Fix a finite loopless multigraph with its individual edge identities, a
palette of $k$ colors, and a distinguished edge $e_0$.  Let
$(T,\varphi_n)$ be a positive-rung extended Tashkinov tree constructed
from an arbitrary certified input coloring, let $\mathcal B$ be a
hierarchy base under $\varphi_n$, and let $C$ be a hierarchy
counterexample over that base.  Write $H=C.\mathit{hierarchy}$, let $q$
be its divider count, and put
\[
 A=\noderef{HierarchyStarredStage}(H,q),\qquad
 \ell=|T|-|A|.
\]
Assume that $\varphi_n$ is proper and leaves exactly $e_0$ uncolored.

Assume, uniformly for every positive-rung maximum ETT on the same graph,
palette, and distinguished edge, with arbitrary certified input and
generating colorings and with rung number smaller than that of $T$, all
six lower-rung assertions of the technical theorem: elementarity;
terminal interchangeability; at every PE slot, unconditional stability
before the support, stability on the preceding ladder tree when the
first-appearance index is greater than one, and disjointness of stored
pairs at distinct PE slots with distinct supports and equal
first-appearance index; the terminal PE component assertion; the
defective-endpoint order assertion; and the mod-coloring assertion.
Assume also that every same-rung maximum ETT with a good hierarchy having
fewer than $q$ dividers is elementary.

Assume that $C$ is lexicographically minimal among hierarchy
counterexamples whose ETTs have the same rung number as $T$ and whose
hierarchies have exactly $q$ dividers, allowing their certified input
coloring, generating coloring, ETT, base, and hierarchy to vary.  Thus no
such $C'$ has
\[
 \noderef{HierarchyCounterexampleLess}
   (\noderef{HierarchyCounterexampleMeasure}(C'),
    \noderef{HierarchyCounterexampleMeasure}(C)).
\]
Finally assume the large-branch estimate
\[
 2n+11\leq |\overline\varphi_n(T_n)|,
\]
where $n$ is the rung number and $T_n$ is the terminal ladder tree.

Let
\[
 b=\min\{|A|+(\ell-1),|T|\},\qquad
 B=\noderef{TreeSequencePrefix}(T,b).
\]
Then
\[
        V(B)\text{ is elementary under }\varphi_n.             \tag{6.6}
\]
All prefixes and equalities here retain the displayed individual edge
identities, including identities of parallel edges.
\end{lemma}
\begin{proof}
We first establish the ordered position of $A$.  If $q>0$, the
positive-index clause of \noderef{HierarchyStarredStage} gives
$A=H.\mathit{stage}(q)$.  The nested and terminal-stage fields of
\noderef{TashkinovHierarchy} therefore give
\[
                     A\prec T.                                 \tag{1}
\]
If $q=0$, the same definition gives $A=\mathcal B.\mathit{core}$,
and `coreInsideFirst' and the terminal-stage field give
\[
                     A\preccurlyeq T.                           \tag{2}
\]
Each relation is a \noderef{TreeSequenceInitialSegment}, so its vertex
and edge equations identify every retained entry, not just the vertex
sets.

We shall also need the following consequence of the hypotheses:
\[
 q=0\quad\Longrightarrow\quad
 V(A)\text{ is elementary under }\varphi_n.                    \tag{3}
\]
Let $i=\operatorname{Fin.last}(n)$ be the last data index and let
$T_n$ be the tree in slot $i$ of the certified series.  The `positive'
field of \noderef{HierarchyCounterexample} gives $n>0$.  Hence the
certificate at $i$ supplied by \noderef{TashkinovSeries} cannot use the
initial constructor of \noderef{TashkinovSeriesStep}; it is exactly one
of PE, SE, and RE.

Suppose first that this constructor is PE.  The reflexive coloring
$\varphi_n$ is stable on $T_n$ relative to its own stored set: both the
vertex-missing-set and controlled-incident-edge clauses in
\noderef{StableColoring} are equalities by reflexivity.  Apply
\noderef{ParallelClosureUnionElementary} to the present ETT, its maximum
property and positivity from $C$, the stated lower-rung package, the PE
kind equation, the assumed properness and unique-uncoloredness, this
reflexive stability, and the present base $\mathcal B$.  Its conclusion
is precisely that the core $A$ is elementary under $\varphi_n$.

It remains to treat SE and RE.  In either case the nonparallel clause of
\noderef{HierarchyBase} identifies
\[
                              A=T_n.                            \tag{4}
\]
Write $\varphi_{n-1}$ for the coloring in the data slot immediately
preceding $i$.  We first prove that $T_n$ is elementary under
$\varphi_{n-1}$ without introducing a rung-zero ETT.  If $n=1$, the
tree and coloring are exactly slots $1$ and $0$ of the original
certified series, so the elementarity conclusion of
\noderef{TechnicalTheoremBase} applies directly.  If $n>1$, apply
\noderef{CertifiedETTRestriction} at target rung $n-1$.  It returns an
actual positive-rung ETT whose displayed tree is exactly the old slot
$T_n$, whose generating coloring is exactly $\varphi_{n-1}$, whose rung
is $n-1<n$, and which has maximum property.  Assertion (i) of the
assumed lower-rung package applied to this returned ETT makes $V(T_n)$
elementary under $\varphi_{n-1}$.  These two cases prove
\[
 V(T_n)\text{ is elementary under }\varphi_{n-1}.              \tag{5}
\]

In the SE case, the constructor contains a
\noderef{MaximumDefectiveChoice}.  Its `selectedStable' field says that
the missing-color set at every vertex of $T_n$ is the same under
$\varphi_{n-1}$ and under its selected coloring.  The SE
`currentColoring' equation in \noderef{TashkinovSeriesStep}, followed by
the generating-coloring equation in
\noderef{ExtendedTashkinovTree}, identifies that selected coloring with
$\varphi_n$.  Therefore (5) transports vertex by vertex to
$\varphi_n$: by \noderef{ColoredVertexSetElementary}, pairwise
disjointness is a statement only about those vertex missing-color sets.
In the RE case, the constructor's `currentColoring'
equation directly identifies the current coloring $\varphi_n$ with
$\varphi_{n-1}$, so (5) already has the desired coloring.  Combining
either conclusion with (4) proves (3).  Notice that the SE transport
uses exactly the certified selected stable coloring, whereas the RE
transport uses exactly the constructor equality; no equality of
generating colorings is asserted in the SE case.

We now compute the prefix length.  From (1)--(2), $|A|\leq|T|$.
Elementary arithmetic with truncated subtraction gives
\[
\begin{array}{rcl}
 \ell=0&\Longrightarrow& b=|T|,\\
 \ell=1&\Longrightarrow& b=|A|,\\
 2\leq\ell&\Longrightarrow& |A|<b<|T|\text{ and }b+1=|T|.
\end{array}                                                     \tag{6}
\]
Indeed, in the first case $|T|=|A|$; in the second
$|T|=|A|+1$; and in the third
$|A|+(|T|-|A|-1)=|T|-1$.  The two defining entrywise equations of
\noderef{TreeSequencePrefix}, together with (1)--(2), have the following
exact consequences.  In the first line of (6), $B=T$; in the second,
$B=A$; and in the third,
\[
                      A\prec B\prec T.                         \tag{7}
\]
For example, $A$ and the prefix of $T$ of length $|A|$ have the same
length and the same vertex and edge entries, which proves their equality
as ordered edge-valued tree-sequences.

We finish by the three cases in (6).

If $\ell=0$, (1) rules out $q>0$, because then $|A|<|T|$.  Hence
$q=0$.  Now $A=B=T$, and (3) proves the desired elementarity of $B$.

Suppose $\ell=1$.  Then $B=A$.  If $q=0$, conclusion (3) again proves
the result.  If $q>0$, apply
\noderef{GoodHierarchyPrefixRepackage} with the old stage index $r=q$.
The cited theorem is applicable because $1\leq q$; it returns an actual
same-rung maximum ETT whose displayed tree is the exact old ordinary
stage $H.\mathit{stage}(q)=A=B$, together with a base and a good
hierarchy having exactly $q-1$ dividers.  Since $q-1<q$, the assumed
fewer-divider elementarity assertion applies to those returned objects
and makes $V(B)$ elementary under $\varphi_n$.

It remains to suppose $2\leq\ell$.  Then (7) holds and $B$ is obtained
from $T$ by deleting exactly its final displayed edge and vertex.  Assume
for contradiction that $V(B)$ is not elementary under $\varphi_n$.
Apply \noderef{GoodHierarchyTerminalPrefixRepackage} to $A\prec
B\preccurlyeq T$.  It supplies an actual ETT $\mathcal E_B$ displayed
on $B$, maximum property, the exact rung equality
$r(\mathcal E_B)=r(\mathcal E)$, a base $\mathcal B_B$, a hierarchy
$H_B$, and a good hierarchy on $H_B$.  It also gives
\[
 q(H_B)=q(H),                                                     \tag{8}
\]
equality of all ordinary stages through index $q$, equality of all
starred stages through that index, and retention of the complete
certified series and all its individual connecting-edge identities.
Together with the assumed non-elementarity, the returned positivity,
goodness, and maximum property package a genuine
\noderef{HierarchyCounterexample} $C_B$ over $\mathcal B_B$.

We next check the first coordinate of the measure, rather than assuming
that truncation preserves it.  Put
\[
       p=\noderef{HierarchyPathNumber}(H),\qquad
       p_B=\noderef{HierarchyPathNumber}(H_B).
\]
The old and new last starred stages agree by the repackaging theorem.
Their tail lengths are respectively $\ell$ and $\ell-1$, by (6).
In the finite candidate set defining $p$, the value $\ell$ is always a
candidate: its endpoint condition concerns an edge position at least
$|A|+\ell=|T|$, of which there is none.  Likewise $\ell-1$ is a
candidate for $p_B$.  Thus both finite minima use their nonempty branch.

If $p=\ell$, minimality of the new candidate set gives
\[
                    p_B\leq\ell-1\leq p.                      \tag{9}
\]
If $p<\ell$, then $p\leq\ell-1$.  We claim that $p$ itself is a
candidate for $H_B$.  Let $s$ be an edge position of $B$ with
$|A|+p\leq s$.  The exact initial-segment relation $B\preccurlyeq T$
maps $s$ to the position with the same numerical value in $T$ and
identifies the individual edge there and both adjacent displayed
vertices.  Since $p$ is an old candidate, that old edge has endpoint
pair equal to those two consecutive vertices.  Transporting this
equation across the three entrywise equalities gives the same endpoint
equation in $B$.  This holds for every such $s$, so $p$ is a new
candidate and $p_B\leq p$.  Together with (9),
\[
                    p_B\leq p.                                \tag{10}
\]
This is exactly the two-branch minimum argument in the definition of
\noderef{HierarchyPathNumber}; it does not confuse the full-tail length
$\ell$ with the path number $p$.

If the inequality in (10) is strict, the first coordinate of
\noderef{HierarchyCounterexampleMeasure} already gives
\[
 \noderef{HierarchyCounterexampleLess}
   (\noderef{HierarchyCounterexampleMeasure}(C_B),
    \noderef{HierarchyCounterexampleMeasure}(C)).              \tag{11}
\]
Suppose instead that $p_B=p$.  The second coordinate must be computed
from the ordinary stage, not from $A$.  Put
\[
                         S=H.\mathit{stage}(q).
\]
The terminal-prefix repackaging theorem identifies the corresponding
ordinary stage of $H_B$ exactly with $S$.  Also
$S\preccurlyeq A\preccurlyeq B$: for $q>0$ the first relation is
equality, while for $q=0$ it is the inclusion of the terminal ladder
tree $H.\mathit{stage}(0)=T_n$ in the hierarchy core.  In the non-PE
case this is the equality supplied by \noderef{HierarchyBase}; in the PE
case it is the left initial-segment field of that base's exact
\noderef{OrderedTreeUnion} certificate.  In particular $|S|\leq|B|$.  Since
$|B|+1=|T|$, subtraction by $|S|$ gives
\[
                 |B|-|S|+1=|T|-|S|,
\]
and hence
\[
                 |B|-|S|<|T|-|S|.                             \tag{12}
\]
These are precisely the new and old second coordinates in
\noderef{HierarchyCounterexampleMeasure}.  This includes $q=0$, where
$S=T_n$ and may be strictly smaller than the PE core $A$.  Equal first
coordinates and (12) again give (11) by
\noderef{HierarchyCounterexampleLess}.

Finally, $C_B$ has the same rung as $C$ and, by (8), the same divider
count.  The rung-local lexicographic minimality hypothesis applied to
$C_B$ forbids (11), a contradiction.  Thus $B$ is elementary in the
case $2\leq\ell$ as well.  The three cases prove (6.6).  The
large-branch estimate is retained as part of the fixed hierarchy
counterexample package needed downstream; this prefix-minimality
argument does not require its numerical value.
\end{proof}
